\documentclass[12pt]{article}
\usepackage[utf8]{inputenc}
\usepackage{amsmath}
\usepackage{amssymb}
\usepackage{graphicx}
\usepackage{minted}
\usepackage{hyperref}
\usepackage{geometry}
\geometry{a4paper, margin=1in}

\title{\textbf{Deconstructing Predictive Coding Networks \\ Part 2: Energy Minimization in Pure PyTorch}}
\author{Sovesh Mohapatra}
\date{\today}

\begin{document}

\maketitle

\section*{Introduction}
In Part 1, we explored the biological plausibility of Predictive Coding Networks (PCNs) and how they flip standard Multi-Layer Perceptrons (MLPs) on their head. We proved that PCNs can replace the global, non-biological \texttt{loss.backward()} of standard Artificial Neural Networks with iterative, local energy minimization.

Today, we will turn that theory into raw PyTorch code. Our goal is to create a network that can learn a simple task without ever calling PyTorch's automatic generalized backpropagation through the whole network!

\section*{The PCN Layer}

A standard PyTorch \texttt{nn.Linear} layer simply maps inputs to outputs. A PCN layer, however, contains dynamic states that require updating during an explicit \textit{inference phase}. Each layer holds a latent ``belief'' and calculates its own local prediction error against the layer below it.

\begin{minted}[frame=lines, fontsize=\small]{python}
import torch
import torch.nn as nn
import torch.nn.functional as F

class PCNLayer(nn.Module):
    def __init__(self, in_features, out_features, activation=torch.tanh):
        super().__init__()
        self.out_features = out_features
        self.activation = activation
        
        # Generative weights predicting the layer below
        self.W = nn.Linear(out_features, in_features, bias=True)
        
        # Internal states for the PCN
        self.latent_state = None
        self.prediction_error = None
        
    def forward(self):
        # Prediction: f(v) * W
        return self.W(self.activation(self.latent_state))
        
    def compute_prediction_error(self, target_state):
        prediction = self.forward()
        self.prediction_error = target_state - prediction
        return self.prediction_error
\end{minted}

Notice we recommend \texttt{torch.tanh}. In standard neural nets, \texttt{ReLU} is the standard. But in iterative PCNs, \texttt{ReLU}'s zero-gradient for negative inputs can cause latent beliefs to "die" during the inference phase, halting learning.

\section*{Phase 1: The Inference Walk}

Unlike feed-forward networks, PCNs don't calculate their output in one step. They enter an \textit{inference phase} where the latent states of all layers recursively undergo gradient descent to minimize total prediction error. The synaptic weights remain entirely frozen during this phase.

\begin{minted}[frame=lines, fontsize=\small]{python}
def inference_step(self, inference_lr):
    total_energy = 0.0
    current_target = self.input_data # Fixed sensory input
    
    # 1. Compute prediction errors bottom-up
    for i in range(self.num_layers):
        error = self.layers[i].compute_prediction_error(current_target)
        # Energy = 1/2 * ||error||^2
        layer_energy = 0.5 * torch.sum(error ** 2)
        total_energy += layer_energy
        current_target = self.layers[i].latent_state
        
    # 2. Extract latent states that require gradients
    latent_states = [layer.latent_state for layer in self.layers 
                     if layer.latent_state.requires_grad]
                     
    if len(latent_states) == 0:
        return total_energy.item()

    # 3. Gradient w.r.t latent states ONLY (Weights are frozen)
    gradients = torch.autograd.grad(total_energy, latent_states)
    
    # 4. Update the latent states (Minimize Energy)
    with torch.no_grad():
        grad_idx = 0
        for layer in self.layers:
            if layer.latent_state.requires_grad:
                layer.latent_state.sub_(inference_lr * gradients[grad_idx])
                grad_idx += 1
                
    return total_energy.item()
\end{minted}

We use PyTorch's \texttt{autograd.grad} explicitly here to descend the energy landscape with respect to the latent nodes. No \texttt{optimizer.step()} is used on the weights.

\section*{Phase 2: The Biological Learning Rule}

Once the inference phase has run for sufficient iterations (e.g., $T=50$) and the latent states have settled into a low-energy configuration, we execute the learning phase. 

This is the biological magic. The generic chain rule is gone. The weight update is executed using \textbf{only} the post-synaptic local prediction error ($\epsilon_{l-1}$) and the pre-synaptic activation ($f(v_l)$).

\begin{minted}[frame=lines, fontsize=\small]{python}
def update_weights(self, learning_rate):
    with torch.no_grad():
        for layer in self.layers:
            error = layer.prediction_error
            pre_synaptic = layer.activation(layer.latent_state)
            
            # Local Hebbian Update Rule
            # delta W = error^T @ pre_synaptic
            delta_W = torch.matmul(error.t(), pre_synaptic) / error.size(0)
            delta_bias = torch.mean(error, dim=0)

            # Update weights locally
            layer.W.weight.add_(learning_rate * delta_W)
            if layer.W.bias is not None:
                layer.W.bias.add_(learning_rate * delta_bias)
\end{minted}

\section*{Next up: The Benchmark Showdown}

We have successfully built a Predictive Coding Network that learns via local energy minimization. But how does this elegant biological algorithm stack up against raw, globally-optimized Backpropagation? 

In \textbf{Part 3}, we will put this PyTorch PCN implementation to the test against a standard identical MLP on nonlinear regression and MNIST image classification to analyze its stability, generative properties, and hardware future.

\vspace{1em}
\noindent \textit{Ready to see the results? Stay tuned for the final data drop!}

\end{document}
